% A cone in 3D
% Author: Marco Daniel
\documentclass{article}
\usepackage{tikz}
%%%<
\usepackage{verbatim}
\usepackage[active,tightpage]{preview}
\PreviewEnvironment{tikzpicture}
\setlength\PreviewBorder{5pt}%
%%%>
\usepackage{tikz-3dplot}
\begin{document}

\usetikzlibrary{calc}

\tdplotsetmaincoords{70}{0}
\begin{tikzpicture}[tdplot_main_coords]
\def\RI{2}
\def\RII{2}

% back for the colored squares
\pgfdeclarelayer{back}
\pgfsetlayers{back,main}

\pgfmathdeclarerandomlist{MyRandomColors}{%
    {red}%
    % {red!25}%
    % {magenta}%
    % {magenta!25}%
    % {olive}%
    % {olive!25}%
    % {brown}%
    % {brown!10}%
    % {violet}%
    % {violet!25}%
    % {gray}%
    % {purple}%
    {yellow}%
    {orange}%
    {orange!60}%
    {cyan}%
    {green}%  
    {blue!60}
    {blue!20}
}%

\draw[opacity=0] (\RI,0)
\foreach \x in {0,60,180} { --  (\x:\RI) 
    node at (\x:\RI) (R1-\x) {} 
    node [pos=0.33] (R1-\x-one) {}
    node [pos=0.66] (R1-\x-two) {} };
\draw[thin] (R1-180.center)
  \foreach \x in {240, 300, 360} { --  (\x:\RI) 
    node at (\x:\RI) (R1-\x) {} 
    node [pos=0.33] (R1-\x-one) {}
    node [pos=0.66] (R1-\x-two) {} };
% \path[fill=gray!30] (\RI,0)
%   \foreach \x in {0,60,120,180,240,300} { --  (\x:\RI)};

\begin{scope}[yshift=2cm]
\draw[ultra thin,fill=blue!10] (\RII,0)
  \foreach \x in {0,60,120,180,240,300,360} { --  (\x:\RII) 
    node at (\x:\RII) (R4-\x) {} 
    node [pos=0.33] (R4-\x-one) {} 
    node [pos=0.66] (R4-\x-two) {} };
\draw[thin] (\RII,0)
  \foreach \x in {0,60,120,180,240,300,360} { --  (\x:\RII) };
\end{scope}

\begin{scope}[yshift=0.66cm]
\draw[ultra thin] (180:\RII)
  \foreach \x in {180,240,300,360} { --  (\x:\RII) node at (\x:\RII) (R2-\x) {}
        node [pos=0.33] (R2-\x-one) {} 
        node [pos=0.66] (R2-\x-two) {} 
  };
\end{scope}
\begin{scope}[yshift=1.33cm]
\draw[ultra thin] (180:\RII)
  \foreach \x in {180,240,300,360} { --  (\x:\RII) node at (\x:\RII) (R3-\x) {}
        node [pos=0.33] (R3-\x-one) {} 
        node [pos=0.66] (R3-\x-two) {} 
  };
\end{scope}

\foreach \x in {0,180,240,300} { \draw (R1-\x.center)--(R4-\x.center); };
% \foreach \x in {60,120} { \draw[dashed] (R1-\x.center)--(R3-\x.center); };
\foreach \x in {240, 300, 360} {
    \draw[ultra thin] (R1-\x-one.center) -- (R4-\x-one.center);
    \draw[ultra thin] (R1-\x-two.center) -- (R4-\x-two.center);
}

\begin{pgfonlayer}{back}
\foreach \x in {180,240,300} {
    \foreach \y in {1,2,3} {

        \pgfmathtruncatemacro\ny{\y + 1}
        \pgfmathtruncatemacro\nx{\x + 60}

        \pgfmathrandomitem{\RandomColor}{MyRandomColors} 
        \path [fill=\RandomColor] (R\y-\x.center) -- (R\ny-\x.center) -- (R\ny-\nx-one.center) -- (R\y-\nx-one.center);
        \pgfmathrandomitem{\RandomColor}{MyRandomColors} 
        \path [fill=\RandomColor] (R\y-\nx-one.center) -- (R\ny-\nx-one.center) -- (R\ny-\nx-two.center) -- (R\y-\nx-two.center);
        \pgfmathrandomitem{\RandomColor}{MyRandomColors} 
        \path [fill=\RandomColor] (R\y-\nx-two.center) -- (R\ny-\nx-two.center) -- (R\ny-\nx.center) -- (R\y-\nx.center);
    }
} 
\end{pgfonlayer}

\end{tikzpicture}

\end{document}